\subsection{From the audited completion theorem to local sections}

\begin{claim}{D.19}{Proved; the section-theoretic consequence of R24}
Let $k/\Q_p$ be finite, and let $X/k$ be smooth, proper,
geometrically connected of genus at least two. For the full arithmetic
\'{e}tale group, the map
\[
 X(k)\longrightarrow
 \{s:G_k\longrightarrow\pi_1(X):s\text{ continuous},\;
                         \mathrm{pr}\circ s=1\}/\Pi_X
\]
is bijective. Here the quotient is by conjugation by the full geometric
group. We derive this from D.18, with no birational-lift hypothesis on $s$.
The argument is the section-theoretic argument of the supplied R24,
whose inputs and valuation comparisons we check explicitly below.
\end{claim}
\begin{proof}
Put $K=k(X)$ and let $\widetilde K\subset K^{\mathrm{sep}}$
be the extension corresponding to $\pi_1(X)$. We first show that,
for every rank-one valuation $w$ of $K$ extending the $p$-adic
valuation, and every chosen prolongation $w^{\mathrm{sep}}$,
restriction induces an isomorphism of profinite groups
\[
 D_{w^{\mathrm{sep}}}(K^{\mathrm{sep}}/K)
       \xrightarrow{\ \sim\ }
 D_{\widetilde w}(\widetilde K/K),
 \qquad \widetilde w=w^{\mathrm{sep}}|_{\widetilde K}.       \tag{D.19.1}
\]
A rank-one ordered value group embeds in $\R$; normalize $w(p)=1$.
For surjectivity, lift an element of the right decomposition group
to $G_K$. Its translate of $w^{\mathrm{sep}}$ is another prolongation
of $\widetilde w$. The kernel $G_{\widetilde K}$ acts transitively
on these prolongations, so an adjustment gives a lift fixing
$w^{\mathrm{sep}}$. The infinite transitivity assertion follows from
the finite normal-extension assertion and compactness of the Galois
group. An automorphism fixing this valuation preserves the normalized
absolute value: two order embeddings of a rank-one value group into
$\R$ differ by a positive scale, fixed by the value of $p$.
It therefore extends continuously to the completion
$E=\widehat{K^{\mathrm{sep}}}$. If it fixes $\widetilde K$, it fixes
its closure $H=\widehat{\widetilde K}\subset E$. D.18 gives the
specified containment $K^{\mathrm{sep}}\subset H$, so that
automorphism is the identity. Decomposition groups are closed:
preservation of membership in a valuation ring is a closed condition
on each finite Galois orbit of algebraic elements. Thus (D.19.1) is
a continuous bijection between compact Hausdorff groups and has a
continuous inverse.

We next prove that
\[
          D_{\widetilde w}(\widetilde K/K)\longrightarrow G_k
                \quad\hbox{has no continuous section}       \tag{D.19.2}
\]
for such a rank-one $w$. Otherwise (D.19.1) would produce a continuous
birational section $\widehat s:G_k\to G_K$ whose image fixes
$w^{\mathrm{sep}}$. The opened Proposition~2.4(2), together with
Definition~1.4 and Observation~1.3, of \citet{D_Koenigsmann2005}
says that, over a finite extension of $\Q_p$, every such birational
section for a smooth complete curve is contained in a decomposition
group of a rational closed-point valuation $\operatorname{ord}_a$,
for some $a\in X(k)$. Choose its prolongation fixed by this same image.

Let $\Sigma=\widehat s(G_k)$ and $M=(K^{\mathrm{sep}})^\Sigma$.
The subgroup is compact and closed, hence $G_M=\Sigma$. The
restrictions $v_1,v_2$ of the two chosen prolongations to $M$ are
henselian: their absolute Galois group fixes one prolongation and
acts transitively on all prolongations, so that prolongation is unique.
They remain nontrivial and rank one. Indeed value groups in an
algebraic valued extension have torsion quotient, which leaves the
chain of convex subgroups unchanged. They are inequivalent because
$v_1(p)=1$ whereas $v_2|_{k^\times}=0$.

We give the contradiction directly, including $p=2$ and nondiscrete
value groups. Choose a nonsquare $d\in k^\times$. It remains a
nonsquare in $M$: a square root would lie in the chosen $\bar k$
and be fixed by $\Sigma$, contradicting the surjection
$\Sigma\to G_k$. Choose a uniformizer $t$ of the local ring at
$a$, and an integer $N>v_1(t)$. Then $z=p^N/t$ satisfies
$v_1(z)>0$ and $v_2(z)<0$. For positive integers $r$ set
\[
 q_r=d+(1-d)\frac{z^r}{1+z^r}.
\]
The denominator is nonzero, since $v_1(z^r)>0$. Direct calculation gives
\begin{align*}
 v_1(q_r/d-1)&=v_1((1-d)/d)+r v_1(z)\longrightarrow+\infty,\\
 v_2(q_r-1)&=v_2(d-1)-r v_2(z)\longrightarrow+\infty.
\end{align*}
For large $r$ these exceed $2v_1(2)$ and $2v_2(2)$, respectively.
Hensel's lemma gives a square root of $q_r/d$ in $(M,v_1)$ and
a square root of $q_r$ in $(M,v_2)$. Their quotient makes $d$ a
square in $M$. This proves (D.19.2); it also supplies the exact
two-valuation argument instead of assuming a general uniqueness slogan.

Now take any continuous section $s:G_k\to\pi_1(X)$.
The opened Pop--Stix valuative localization theorem
\citep[Theorem~26 in the author preprint]{D_PopStix2017}
applies to a smooth hyperbolic geometrically connected curve over
a finite extension of $\Q_p$, with no assumed point or birational
lift. It supplies a valuation $w$ of $K$ extending the $p$-adic
valuation and a prolongation $\widetilde w$ to $\widetilde K$ such that
\[
                    s(G_k)\subset D_{\widetilde w}.
\]
The rational rank of $\Gamma_w/\Gamma_k$ is at most one.
For clarity, two values rationally independent modulo $\Gamma_k$
would be the values of elements $u,v\in K^\times$; distinct monomials
in any polynomial in $u,v$ over $k$ would have distinct valuations,
so no such polynomial could vanish. This would give transcendence
degree at least two, contrary to $K=k(X)$. Since $\Gamma_k$ has
rational rank one, the ordered rank of $w$ is at most two.

Rank one is excluded by (D.19.2). In rank two let $u$ be the
rank-one coarsening. The nonzero proper convex subgroup of
$\Gamma_w$ extends uniquely through an algebraic value-group
extension, by taking all elements a positive multiple of which
lies in that subgroup. Consequently $s(G_k)$ preserves the induced
prolongation of $u$. If $u$ is nontrivial on $k$, its restriction
is equivalent to the $p$-adic valuation, since that is the restriction
of $w$ and a nontrivial coarsening of a rank-one valuation is itself.
Normalization and (D.19.2) again exclude it.

Thus $u$ is trivial on $k$. Properness gives a center on $X$.
It is not the generic point because $u$ is nontrivial, hence is a
closed point $a$. Its valuation ring dominates the discrete valuation
ring $\mathcal O_{X,a}$. Writing each element of $K^\times$ as a
power of a uniformizer times a unit of that ring shows that $u$
is equivalent to $\operatorname{ord}_a$. The refinement $w$ is
the composition with the unique $p$-adic valuation on the finite
extension $k(a)/k$. A chosen horizontal prolongation selects a
geometric point over $a$, so its decomposition group maps to the
corresponding conjugate of $G_{k(a)}\subset G_k$. Since $s$ maps
onto $G_k$, this inclusion forces $k(a)=k$.

In the full \'{e}tale tower, the horizontal inertia at $a$ is
trivial. Its decomposition group maps isomorphically to $G_k$
and is precisely the point section, as is also recorded in
\citet[Section~6.1]{D_PopStix2017}. Therefore $s$ has the class
of $s_a$. Changes of geometric lift conjugate by $\Pi_X$,
which is exactly the ambiguity in the claim. This proves surjectivity.

For injectivity suppose $s_a$ and $s_b$ are geometrically conjugate.
Use $a$ to define $i_a:X\to J=\operatorname{Jac}(X)$ by
$x\mapsto[x-a]$. For every integer $n>0$, the map $[n]:J\to J$
is finite \'{e}tale, including $p\mid n$, because the field has
characteristic zero. The fibre at $i_a(a)=0$ has a fixed lift.
Conjugacy of sections preserves the existence of such a lift in
every finite \'{e}tale cover, hence
\[
                     i_a(b)\in\bigcap_{n\ge1}nJ(k).
\]
The compact $p$-adic analytic group $J(k)$ has a torsion-free open
subgroup $U\simeq\Z_p^{[k:\Q_p]\dim J}$ by the local logarithm.
For $q=[J(k):U]$ we have $qJ(k)\subset U$, so the displayed
intersection lies in $\bigcap_r p^rU=0$. Thus $[b-a]=0$.
If $a\ne b$, a rational function with divisor $b-a$ would have
degree one as a map to $\mathbb P^1$, contradicting positive genus.
Thus $a=b$ and the proof is complete.
\end{proof}

The argument verifies the local proper assertion of R24 through its
specified completion theorem. It supplies no almost-everywhere
integrality or global finite-descent equality over a number field;
the continuous global obligations D1--D4 remain separate.
